\section*{Supplemental information}

%%% Supplemental items are supplied as a separate PDF (supplement.tex);
%%% only titles and legends belong in the main document.
\setcounter{figure}{0}\renewcommand{\thefigure}{S\arabic{figure}}
\setcounter{table}{0}\renewcommand{\thetable}{S\arabic{table}}

\refstepcounter{figure}\label{fig:length_compliance}
\refstepcounter{figure}\label{fig:supplementary_overview}
\refstepcounter{table}\label{tab:training_cutoff_overlap}

\begin{itemize}
	\item Figure S1. \textbf{Summary length compliance per method:} Proportion of generated summaries that fell within the target length of 15--100 words, were too short, or were too long, shown for each of the 62 evaluated methods. Methods are ordered by decreasing share of within-bounds summaries, with the within-bounds percentage annotated at the right of each bar. \emph{Related to STAR Methods.}
	\item Figure S2. \textbf{Model performance across metrics:} Overview of the performance of all evaluated models across all metrics. \emph{Related to Figure 3.}
	\item Table S1. \textbf{Training cutoff dates and potential overlap for all evaluated models:} 
The training cutoff indicates the latest point in time up to which a model may have been trained; the source of this information is given in a separate column (confirmed, community, estimated, or unknown). Potential overlap describes the percentage of benchmark articles published on or before the respective cutoff date and therefore represents the theoretical maximum share of articles that could have been present in the model's training data. Closed-source models are marked accordingly. Models are ordered by decreasing potential overlap. \emph{Related to Table 3.}
\end{itemize}
